\documentclass[10pt,a4paper]{amsart}
\usepackage[margin=19mm]{geometry}
\usepackage{amsmath,amssymb,mathtools}
\usepackage[scr=rsfs,cal=euler]{mathalfa}
\usepackage{xfrac}
\usepackage[unicode,hidelinks,bookmarks=false]{hyperref}
\newcommand{\ie}{i.e.\ }
\newcommand{\cf}{cf.\ }
\newcommand{\resp}{respectively\ }
\newcommand{\Subsection}{\subsection}
\providecommand{\nobreakdash}{\nobreak\hbox{-}}
\setlength{\emergencystretch}{2em}
\multlinegap0pt
\usepackage{bm}
\usepackage{bbm}
\usepackage{dsfont}
\usepackage{xspace}
\usepackage{cancel}
\usepackage{mathtools}
\usepackage{amsthm}
\newtheorem{theorem}{Theorem}
\newcommand{\lavg}{\ell_{\mathrm{avg}}}
\title[Lean-Certified Infinite Counterexamples to Written on the Wa]{Lean-Certified Infinite Counterexamples to Written on the Wall II Conjecture 194 (Theorem 1)}
\author{Cameron Beeley}
\date{Source: arXiv:2609.19195v1 (CC BY 4.0)}
\begin{document}
\maketitle

\noindent\emph{Source and licence.} Lean-Certified Infinite Counterexamples to Written on the Wall II Conjecture 194. Version arXiv:2609.19195v1 (CC BY 4.0), distributed under the Creative Commons Attribution 4.0 International licence, as recorded in the arXiv licence metadata for this exact version and shown on the abstract page https://arxiv.org/abs/2609.19195v1. Full rights notice: licenses/arxiv-2609.19195-CC-BY-4.0.txt.\par\smallskip

\noindent\textit{Editorial scope.} Counted unit: the Theorem printed as Theorem 1 of the article (source0.tex lines 117--128, source label \texttt{thm:family}), delivered together with the article's own complete original proof of it (source0.tex lines 130--172, one \texttt{proof} environment of 43 lines). No mathematics is written, repaired or re-derived: statement and proof are the article's own text line for line. Required context retained uncounted: none. Every definition, display and label of those lines is retained unchanged. Omitted from this package and not redistributed: 1--116, 129--129, 173--404. Nothing inside a retained excerpt is omitted or altered: every retained body is delivered exactly as the article's lines are written, so no omission receipt of any kind accompanies this package. Citations: the retained excerpts carry no citation command at all, so no bibliography entry of the article is transcribed and no reference is redistributed. Reference adaptation: none was needed and none was made; every reference inside the delivered unit resolves inside it, so no unresolved reference or citation is produced. Printed slips kept literal: no printed word, symbol, hypothesis or number of the counted statement or of its proof was altered. Layout and class: the article's own class is \texttt{article} with packages including fontenc, lmodern, amsmath, amssymb, amsthm, geometry, booktabs, hyperref, cleveref; the wrapper is the lane's pinned \texttt{amsart} editorial wrapper at 10pt on A4, which restates the theorem-like environments the retained text needs and adds the article's own macro definitions that the retained text needs (\texttt{\textbackslash lavg}), with extra wrapper packages (bm, bbm, dsfont, xspace, cancel, mathtools). The delivered numbering is the unit's own. Assets: no retained span contains a figure environment, a graphics-inclusion command, a PDF-page inclusion, a table or a TikZ picture; no image file is copied into this package, the package declares no asset dependency and no image file is redistributed with it. Licence: version arXiv:2609.19195v1 of this article is distributed under the Creative Commons Attribution 4.0 International licence, as recorded in the arXiv licence metadata for this exact version and shown on the abstract page https://arxiv.org/abs/2609.19195v1. A full rights notice is delivered at licenses/arxiv-2609.19195-CC-BY-4.0.txt. Characters: every retained excerpt is pure ASCII, so the delivered engine, XeLaTeX with its default Latin Modern text font, reports no missing character.
\begin{theorem}\label{thm:family}
Let $s\geq 1$ and $3\leq t<r$, and put
\[
  M_s(r,t)=r^2-2r+\bigl(s(r-2)-1\bigr)t.
\]
For every $m\geq M_s(r,t)$, the connected graph $F^{(s)}_{m,r,t}$ satisfies
\[
  \alpha(F^{(s)}_{m,r,t})\leq 1+\lavg(F^{(s)}_{m,r,t})
\]
but has no Hamiltonian path.  Equality holds in the displayed inequality if and only if
$m=M_s(r,t)$.
\end{theorem}

\begin{proof}
The graph is connected because $C$ is nonempty and is completely joined to $H$.
We first determine its independence number.  An independent set meeting $C$ contains
at most one vertex of $C$, no vertex of $H$, and at most one vertex of each $B_i$; it
therefore has size at most $t+1\leq r$.  An independent set disjoint from $C$ contains at
most one vertex from each clique $B_i\cup\{h_i\}$, together with at most the $r-t$
remaining vertices of $H$.  Its size is therefore at most $r$, and this bound is attained
by $H$.  Hence
\[
  \alpha(F^{(s)}_{m,r,t})=r.
\]

The local independence numbers have four forms.  If $c\in C$, then
$N(c)=(C\setminus\{c\})\cup H$, so $\ell(c)=r$.  If $1\leq i\leq t$, then
$N(h_i)=C\cup B_i$ is the disjoint union of two cliques, so $\ell(h_i)=2$.  Each of the
other $r-t$ vertices of $H$ has neighbourhood $C$ and local independence number $1$.
Finally, if $b\in B_i$, then $N(b)=(B_i\setminus\{b\})\cup\{h_i\}$ is a clique, so
$\ell(b)=1$.  Consequently
\begin{equation}\label{eq:localsum}
  \sum_{v\in V(F^{(s)}_{m,r,t})}\ell(v)
  =mr+2t+(r-t)+st
  =mr+r+(s+1)t.
\end{equation}
Since $|V(F^{(s)}_{m,r,t})|=m+r+st$, the conjectured hypothesis is equivalent to
\begin{align*}
  r
  &\leq 1+\frac{mr+r+(s+1)t}{m+r+st}\\
  &\Longleftrightarrow
  (r-1)(m+r+st)\leq mr+r+(s+1)t\\
  &\Longleftrightarrow
  m\geq r^2-2r+\bigl(s(r-2)-1\bigr)t.
\end{align*}
This also shows that equality holds exactly when $m=M_s(r,t)$.

Finally, fix $i\leq t$.  The only edges from the nonempty set $B_i$ to its complement are
incident with $h_i$.  A spanning path must use at least one such edge.  It cannot use two:
both would exhaust the two path edges available at $h_i$, leaving no path edge from
$B_i\cup\{h_i\}$ to the nonempty remainder of the graph.  Thus exactly one edge of a
spanning path would cross from $B_i$ to its complement, forcing one endpoint of the path
to lie in $B_i$.  The pairwise disjoint sets $B_1,\dots,B_t$ would therefore require at
least $t\geq 3$ distinct endpoints, whereas a path has only two.  Hence
$F^{(s)}_{m,r,t}$ is not traceable.
\end{proof}

\end{document}
